\section{The ES coefficient boundary with every signed state retained}
\label{sec:es-total-extension}
The starting coefficient equation and rank-three inclusion were proposed
in the received defining-prime account \cite[eqs.\ (3)--(8)]{DefiningPrimeReceived}.
We prove them here and extend them to the entire signed quartic algebra,
its two original conductor receivers, and every inverse exterior norm.
The original four labels and polynomial provenance remain
\cite{ES488,Tao}; FS1--25 and GD14 specify the complete signed algebra
and the unchanged physical return. Integral closure below is a specified
finite morphism of rings, not a rescaling or omission of their coordinates.

\subsection{The exact coefficient inclusion and its conductor ideal}
Put
\[
 R=\mathbb C[A,A^{-1},B,C],\qquad
 \mathscr S=R[D]/(625AD^3-125CD^2+25BC^2D-4C^4).
 \tag{ET1}
\]
Here $A$ is the original quartic coefficient, not the amplification
$a_{\rm amp}$. The retained distinguished prime parameter is $P$.
Define
\[
 \mathscr N=R[P]/\left(P^3+A^{-1}P^2+(B/A)P+\frac{4C}{5A}\right),
 \quad D=-CP/5,\quad W=A^{-1}(B+P+AP^2).
 \tag{ET2}
\]
Both displayed polynomials become monic by multiplying by units.
Consequently $\mathscr S$ and $\mathscr N$ are free over $R$ in the
ordered bases $(1,D,D^2)$ and $(1,P,P^2)$, respectively. Substitution
in ET1 gives zero. The inclusion matrix and its determinant are
\[
 K(C)=\diag(1,-C/5,C^2/25),\qquad \det K=-C^3/125.
 \tag{ET3}
\]
This diagonal map is injective over the domain $R$; after inverting
$C$ it is an isomorphism. Solving ET2 for $B,C,D$ gives
\[
 B=AW-P-AP^2,\quad C=-5APW/4,\quad D=AP^2W/4,
 \qquad \mathscr N=\mathbb C[A,A^{-1},P,W].
 \tag{ET4}
\]
The inverse ring maps follow by substituting these three expressions
in ET2. Thus $\mathscr N$ is a domain and is integrally closed:
a polynomial ring over a field is a unique factorization domain,
localization preserves that property, and a unique factorization
domain is integrally closed by applying prime factorization to the
denominator of an element satisfying a monic equation.
Since the rings have the same fraction field and $\mathscr N$ is
finite integral over $\mathscr S$, it is the full integral closure.
Indeed every fraction integral over $\mathscr S$ is integral over
$\mathscr N$ by the same monic equation, hence belongs to $\mathscr N$.

The quotient and the two expressions of the conductor ideal are
\[
 \mathscr N/\mathscr S=R/(C)\oplus R/(C^2),\qquad
 (\mathscr S:\mathscr N)=C^2\mathscr N=(C,D)^2\mathscr S .
 \tag{ET5}
\]
Here $(\mathscr S:\mathscr N)=\{x\in\mathscr N:x\mathscr N\subset\mathscr S\}$.
For a proof retaining every coefficient, write an element of $\mathscr S$
as $a+CbP+C^2cP^2$; the harmless constant units $-1/5,1/25$
are retained in ET3, and do not change divisibility by $C$.
Set $a_0=A^{-1},b_0=B/A,c_0=4C/(5A)$.
Its product with $P$ has coefficients
$(-C^2cc_0,\ a-C^2cb_0,\ Cb-C^2ca_0)$.
Membership in $\mathscr S$ forces $C\mid a$ and $C\mid b$.
In its product with $P^2$, the $P^2$ coefficient is
$a-Cba_0+C^2c(a_0^2-b_0)$, so $C^2\mid a$.
These conditions say exactly that the original element is in $C^2\mathscr N$.
Conversely $C^2\mathscr N$ is an ideal of $\mathscr N$ contained in
$\mathscr S$. Its $R$ generators $C^2,C^2P,C^2P^2$ are
$C^2,-5CD,25D^2$, proving the last equality.

\subsection{The complete differential and its signed fibres}
The coefficient map $\nu:(A,P,W)\mapsto(A,B,C,D)$ of ET4 has Jacobian
\[
 J_\nu=\begin{pmatrix}
 1&0&0\\
 W-P^2&-1-2AP&A\\
 -5PW/4&-5AW/4&-5AP/4\\
 P^2W/4&APW/2&AP^2/4
 \end{pmatrix}.
 \tag{ET6}
\]
In row orders $ABC,ABD,ACD,BCD$, its four maximal minors are
\[
 \left(
 \frac{5A}{4}(2AP^2+AW+P),\
 -\frac{AP}{4}(2AP^2+2AW+P),\
 \frac{5A^2P^2W}{16},\
 -\frac{5A^2P^4W}{16}
 \right).
 \tag{ET7}
\]
Expansion of each two-by-two minor after the first row gives the first
three; expanding the remaining determinant gives the fourth.
Their common zero locus is exactly
\[
 L_0=\{P=W=0\},\qquad
 L_1=\{W=0,\ P=-1/(2A)\}.
 \tag{ET8}
\]
Indeed the third minor forces $P=0$ or $W=0$.
If $P=0$, the first forces $W=0$; if $W=0$, the first two force
$P(1+2AP)=0$. The first row and the nonzero $W$ derivative of $B$
show that the rank is exactly two at both loci.
In coordinate order $(dA,dP,dW)$ their kernels are, respectively,
$\mathbb C(0,A,1)$ and $\mathbb C(0,1,0)$.

These loci have explicit fibres in the original signed cover
$\mathbb C[r,\sigma]/(h,\sigma^2-h')$, where
$h=Ar^4+r^3+Br^2+Cr+D$. They are
\[
 \begin{array}{ll}
 L_0:&h=r^3(Ar+1),\\[2pt]
 L_1:&h=Ar^2(r+1/(2A))^2 .
 \end{array}
 \tag{ET9}
\]
At $L_0$, reduction in the root-zero primary gives
$\mathbb C[r,\sigma]/(r^3,\sigma^2-3r^2)$, of length six.
The other root is $-1/A$, with $h'=-1/A^2$, hence its two
distinct signs are $\sigma=\pm i/A$. At $L_1$, each primary
has unchanged local coordinate $\epsilon=r-r_0$ and algebra
$\mathbb C[\epsilon,\sigma]/(\epsilon^2,\sigma^2-\epsilon/(2A))$,
of length four, for $r_0=0,-1/(2A)$. These follow by expanding $h'$
and reducing modulo the displayed primary powers, not by discarding
the cofactor. Distinct primary factors are separated by Bezout
polynomials since their root differences are units. Dimensions add
to eight in both cases.

The original coefficient receivers give a further exact differential
square:
\[
 \nu_U=\Phi\nu,\quad \nu_W=\Psi\nu,\quad
 T_A\nu_U=\nu_W,\quad T_A\,d\nu_U=d\nu_W .
 \tag{ET10}
\]
Both $\Phi,\Psi$ are the fixed isomorphisms of FC, so the ranks and
kernels in ET8 remain exactly those of the transported differentials.
This is a calculation of $d\nu$, not a replacement of the nonzero
conductor determinant FC23.

\subsection{A rank-twenty-four comparison retaining the whole cover}
Define over the actual coefficient ring $\mathscr S$
\[
 \mathscr E=\mathscr S[r,\sigma]/
 \left(r^4+A^{-1}r^3+(B/A)r^2+(C/A)r+D/A,\
       \sigma^2-(4Ar^3+3r^2+2Br+C)\right).
 \tag{ET11}
\]
Successive division by the two monic polynomials proves unique
remainders in the full ordered basis
\[
 e=(1,r,r^2,r^3,\sigma,\sigma r,\sigma r^2,\sigma r^3).
 \tag{ET12}
\]
It is therefore free of rank eight over $\mathscr S$.
Set $\mathscr E_{\mathscr N}=\mathscr N\otimes_{\mathscr S}\mathscr E$.
Both algebras are free of rank twenty-four over $R$, with ordered
bases $(e,De,D^2e)$ and $(e,Pe,P^2e)$. Their exact inclusion is
\[
 \mathcal K=K(C)\otimes I_8,\quad
 \det\mathcal K=C^{24}/5^{24},\quad
 \mathscr E_{\mathscr N}/\mathscr E
 =(R/(C))^8\oplus(R/(C^2))^8 .
 \tag{ET13}
\]
This follows coefficient by coefficient in ET12.
No choice of a root or a sign occurs in this map. It base-changes the
entire signed cover along the distinguished-parameter inclusion.

The exact algebra conductor of this extension is
\[
 (\mathscr E:\mathscr E_{\mathscr N})=C^2\mathscr E_{\mathscr N}.
 \tag{ET14}
\]
To prove necessity, write $x=\sum_{j=0}^7 n_j e_j$ with $n_j\in\mathscr N$.
If $x\mathscr E_{\mathscr N}\subset\mathscr E$, multiplication by every
scalar $n\in\mathscr N$ and uniqueness in ET12 give
$n_j\mathscr N\subset\mathscr S$ for each $j$. ET5 gives
$n_j\in C^2\mathscr N$. Conversely every coefficient of
$C^2\mathscr E_{\mathscr N}$ lies in $\mathscr S$, and this is an
ideal of $\mathscr E_{\mathscr N}$, so it is contained in the conductor.
ET14 describes this particular extension; it does not assert that
$\mathscr E_{\mathscr N}$ is the total integral closure of the signed
cover at its own ramification.

At $C=0$, $\mathcal K$ has rank eight and kernel the sixteen-dimensional
span of $(De,D^2e)$. Tensoring the quotient exact sequence with $R/(C)$
gives the full exact sequence
\[
 0\longrightarrow (R/(C))^{16}\longrightarrow
 \mathscr E/C\mathscr E\longrightarrow
 \mathscr E_{\mathscr N}/C\mathscr E_{\mathscr N}
 \longrightarrow(R/(C))^{16}\longrightarrow0 .
 \tag{ET15}
\]
Here the left term is $\operatorname{Tor}_1^R(
 \mathscr E_{\mathscr N}/\mathscr E,R/(C))$.
For each summand $R/(C^j)$, the free resolution
$0\to R\xrightarrow{C^j}R\to R/(C^j)\to0$
computes both its specialization and Tor term as $R/(C)$.
Thus the left injection retains all sixteen killed directions;
they have not been identified with absent root labels.

\subsection{All root and sign actions, and the distinguished-parameter pole}
In the increasing-power root basis define
\[
 X(D)=\begin{pmatrix}
 0&0&0&-D/A\\1&0&0&-C/A\\0&1&0&-B/A\\0&0&1&-1/A
 \end{pmatrix},\quad
 Z(D)=4AX(D)^3+3X(D)^2+2BX(D)+CI_4.
 \tag{ET16}
\]
Every entry of the latter matrix is
\[
 Z(D)=\begin{pmatrix}
 C&-4D&D/A&D(2AB-1)/A^2\\
 2B&-3C&-(4AD-C)/A&(2ABC+AD-C)/A^2\\
 3&-2B&-(3AC-B)/A&-(4A^2D-2AB^2-AC+B)/A^2\\
 4A&-1&-(2AB-1)/A&-(3A^2C-3AB+1)/A^2
 \end{pmatrix}.
 \tag{ET17}
\]
The complete actions on ET12 are
\[
 X_8(D)=\diag(X(D),X(D)),\qquad
 Y_8(D)=\begin{pmatrix}0&Z(D)\\I_4&0\end{pmatrix}.
 \tag{ET18}
\]
They commute, satisfy $h(X_8)=0$ and
$Y_8^2=4AX_8^3+3X_8^2+2BX_8+CI_8$.
These identities follow from monic reduction for $X$ and direct block
multiplication for $Y_8$.

The two scalar matrices for $D$ are
\[
 M_D=\begin{pmatrix}
 0&0&4C^4/(625A)\\1&0&-BC^2/(25A)\\0&1&C/(5A)
 \end{pmatrix},\quad
 M_P^{\mathscr N}=\begin{pmatrix}
 0&0&-4C/(5A)\\1&0&-B/A\\0&1&-1/A
 \end{pmatrix},\quad M_D^{\mathscr N}=-CM_P^{\mathscr N}/5.
 \tag{ET19}
\]
For $Q=X_8,Y_8$, let $Q_0=Q(0),Q_1=\partial_DQ$.
ET16--18 show that $Q=Q_0+DQ_1$ exactly.
Thus all twenty-four dimensional root and sign actions are
\[
 \mathcal Q_{\mathscr S}=I_3\otimes Q_0+M_D\otimes Q_1,\quad
 \mathcal Q_{\mathscr N}=I_3\otimes Q_0+M_D^{\mathscr N}\otimes Q_1,
 \quad \mathcal K\mathcal Q_{\mathscr S}
 =\mathcal Q_{\mathscr N}\mathcal K .
 \tag{ET20}
\]
The last identity follows from $KM_D=M_D^{\mathscr N}K$,
which is checked entrywise using ET3 and ET19.
It is also the scalar multiplication identity in the proved inclusion.

On $C\ne0$, multiplication by the retained $P=-5D/C$ in the
original coefficient basis has the exact matrix and inverse
\[
 M=K^{-1}M_P^{\mathscr N}K
 =\begin{pmatrix}
 0&0&-4C^3/(125A)\\-5/C&0&BC/(5A)\\0&-5/C&-1/A
 \end{pmatrix},
 \qquad
 M^{-1}=\begin{pmatrix}
 -5B/(4C)&-C/5&0\\25/(4C^2)&0&-C/5\\-125A/(4C^3)&0&0
 \end{pmatrix}.
 \tag{ET21}
\]
Multiplication in both orders gives $I_3$.
The characteristic polynomial is the unchanged cubic ET2;
$\det M=-4C/(5A)$ and $\det M^{-1}=-5A/(4C)$.
The full signed-cover action is $\mathcal M=M\otimes I_8$.
It commutes with both $\mathcal X_{\mathscr S}$ and
$\mathcal Y_{\mathscr S}$ in ET20: $M=-5M_D/C$ commutes with
$M_D$, and scalar and root/sign tensor factors commute.
Thus the original signed action, its scalar extension and its inverse
are specified on the same twenty-four dimensional object.

Fix $A\ne0$ and $B$ and let complex $C$ tend to zero.
Retaining its phase, the exact residue and inverse leading term are
\[
 C M\longrightarrow N=-5(E_{10}+E_{21}),\quad
 \rank(N,N^2,N^3)=(2,1,0),\qquad
 C^3M^{-1}\longrightarrow-\frac{125A}{4}E_{20}.
 \tag{ET22}
\]
Indices are zero-based. For $\mathcal M$ these ranks are
$(16,8,0)$ and the leading inverse rank is eight.
The eight chains are $e_j\mapsto-5De_j\mapsto25D^2e_j\mapsto0$.
This is the coefficient lattice residue; it does not change the
characteristic cubic or replace the actual root and sign operators.

\subsection{The full original conductor and metric return}
Let $F_U,F_W$ be exactly GD14, with the original shift $\kappa$,
all FC24 inverse moments and the original Gamma masses.
For a column in three consecutive copies of ET12, define
\[
 \mathfrak F_U=I_3\otimes F_U,\qquad
 \mathfrak F_W=I_3\otimes F_W,\qquad
 \mathfrak T=T_A^{\oplus6},\qquad
 \mathfrak T\mathfrak F_U=\mathfrak F_W .
 \tag{ET23}
\]
Each pair in one copy returns $R(r)+\sigma S(r)$ as
$(T_A^{-1}R(S'-\kappa),T_A^{-1}S(S'-\kappa))$ at the source
and $(R(S'-\kappa),S(S'-\kappa))$ at the target.
These substitutions and FC24 give the inverses of both receivers.
No moments vary with $A,B,C,D$.

For every one of the exact maps $\mathcal Z$ in ET13, ET20 or ET21,
its two physical versions are
\[
 \mathcal Z_U=\mathfrak F_U\mathcal Z\mathfrak F_U^{-1},\quad
 \mathcal Z_W=\mathfrak F_W\mathcal Z\mathfrak F_W^{-1},\quad
 \mathfrak T\mathcal Z_U=\mathcal Z_W\mathfrak T .
 \tag{ET24}
\]
For an inclusion, source and target copies in this formula carry their
specified $\mathscr S$ and $\mathscr N$ bases. Multiplying ET23 proves
the equality in either interpretation. This supplies the exact
morphism to the fixed conductor for the entire cover.

Write $G_{U,r},G_{W,r}$ for the full degree-three coefficient Grams
of the two single-component receivers in GD14, with all original
moments, centers and Gamma masses as evaluated in RS21--26 and MR1--7.
The metrics on this actual direct sum are, respectively,
\[
 \Gamma_U^{(24)}=I_3\otimes\diag(G_{U,r},G_{U,r}),\qquad
 \Gamma_W^{(24)}=I_3\otimes\diag(G_{W,r},G_{W,r}).
 \tag{ET25}
\]
This follows by adding the squared norms of the six actual components.
In particular $G_{U,r}$ and $G_{W,r}$ are not equated.
For either positive inner block $\Gamma$, the metric square of
$M^{-1}\otimes I_8$ is
\[
 (I_3\otimes\Gamma)^{-1}
 (M^{-1}\otimes I_8)^*(I_3\otimes\Gamma)
 (M^{-1}\otimes I_8)
 =((M^{-1})^*M^{-1})\otimes I_8 .
 \tag{ET26}
\]
Every original metric entry cancels by this exact tensor identity;
no metric was replaced by a convenient one. Thus its three Euclidean
singular values occur with multiplicity eight in both physical spaces.

Here is a fully evaluated finite-parameter formula. For $C\ne0$ put
\[
 x=-5B/(4C),\quad y=-C/5,\quad z=25/(4C^2),\quad w=-125A/(4C^3).
\]
Then $(M^{-1})^*M^{-1}$ is
\[
 Q=\begin{pmatrix}
 |x|^2+|z|^2+|w|^2&\bar x y&\bar z y\\
 \bar y x&|y|^2&0\\
 \bar y z&0&|y|^2
 \end{pmatrix}.
 \tag{ET27}
\]
Its positive eigenvalues $\lambda_1\ge\lambda_2\ge\lambda_3>0$
are exactly the three roots, with their multiplicities, of
\[
 \lambda^3-a_1\lambda^2+a_2\lambda-a_3=0,\qquad
 \begin{aligned}
 a_1&=|x|^2+|z|^2+|w|^2+2|y|^2,\\
 a_2&=|y|^2(|x|^2+|z|^2+2|w|^2)+|y|^4,\\
 a_3&=|y|^4|w|^2=25|A|^2/(16|C|^2).
 \end{aligned}
 \tag{ET28}
\]
Indeed expansion of the determinant of $Q-\lambda I$ gives precisely
these coefficients; positivity follows since $M^{-1}$ is invertible.
The roots can all be evaluated by one quadratic. Set
$t=|y|^2$, $r=|x|^2+|z|^2+|w|^2$, $q=|w|^2$.
Expansion, or substitution in ET28, factors the cubic as
$(\lambda-t)(\lambda^2-(r+t)\lambda+tq)$. Hence
\[
 \lambda_1=\frac{r+t+\sqrt{(r+t)^2-4tq}}2,\quad
 \lambda_2=t,\quad
 \lambda_3=\frac{2tq}{r+t+\sqrt{(r+t)^2-4tq}} .
 \tag{ET28a}
\]
The quadratic is positive at zero and negative at $t$, since
$q-r=-|x|^2-|z|^2<0$. Thus its roots satisfy
$\lambda_1>t>\lambda_3>0$; this also proves that its discriminant
is strictly positive. The rationalized last expression retains every
coefficient and avoids cancellation. In particular eight inverse
singular values are exactly $|C|/5$, and eight more merely approach
that scale. No root or metric constant remains unnamed.
Let $s_i=\sqrt{\lambda_i}$. For $1\le j\le24$, set
\[
 n_1=\min(j,8),\quad n_2=\min(\max(j-8,0),8),\quad
 n_3=\max(j-16,0).
\]
The exact original-metric exterior inverse norm is
\[
 \left\|\bigwedge^j\mathcal M_U^{-1}\right\|
 =\left\|\bigwedge^j\mathcal M_W^{-1}\right\|
 =s_1^{n_1}s_2^{n_2}s_3^{n_3}.
 \tag{ET29}
\]
A singular-value decomposition in orthonormal bases proves that
exterior singular values are the corresponding products, and ET26
supplies the multiplicity eight. Thus ET29 covers every exterior degree,
including the full determinant
$\|\bigwedge^{24}\mathcal M^{-1}\|=(5|A|/(4|C|))^8$.

There are two equal positive singular values $5$ in the residue $N$.
Continuity of the metric square applied to $CM$ shows that the two
largest singular values of $M$ are $5/|C|(1+o(1))$.
Using the exact determinant in ET21 gives the third one as
$4|C|^3/(125|A|)(1+o(1))$. Inverting these three values proves
\[
 s_1=\frac{125|A|}{4|C|^3}(1+o(1)),\qquad
 s_2=\frac{|C|}{5}(1+o(1)),\qquad
 s_3=\frac{|C|}{5}(1+o(1)).
 \tag{ET30}
\]
This also follows for the largest inverse value directly from ET22.
For $d_A=125|A|/4$ and all twenty-four exterior degrees, ET29 gives
\[
 \left\|\bigwedge^j\mathcal M^{-1}\right\|\sim
 \begin{cases}
 d_A^j|C|^{-3j},&1\le j\le8,\\
 d_A^8\,5^{-(j-8)}|C|^{j-32},&8\le j\le24.
 \end{cases}
 \tag{ET31}
\]
The formulas agree at $j=8$; at $j=24$ the constant equals
$(5|A|/4)^8$, checking the complete determinant.
They hold for every complex approach with fixed $A\ne0,B$.
Equations ET21--31 establish an actual cubic inverse pole and all its
complementary directions in the stated original receiving norms.
The unchanged linear conductor is still the explicit FC24 isomorphism;
the pole belongs to the displayed distinguished-parameter action.

For comparison, the full coefficient inclusion itself has inverse
singular values $1,5/|C|,25/|C|^2$, each with multiplicity eight in
ET25, sorted according to their actual sizes. This follows from ET3
and the same tensor metric identity. Its determinant inverse norm is
$5^{24}/|C|^{24}$, consistent with ET13. The inclusion pole, the
distinguished-parameter action pole and the fixed conductor have thus
been related by exact maps while retaining their different matrices.

\subsection{The inverse-exterior estimate with the actual conductor included}
Let $\gamma_1\ge\cdots\ge\gamma_8>0$ be the singular values of
$(T_A\oplus T_A)^{-1}$ in its original target-to-source metrics.
They are the four singular values of the exact FC24 inverse
$A_{3,s}^{-1}$, each repeated twice, using precisely its original
$\rho_j^{(s)},h_j,\mu_v$ and centers. In particular
$\prod_{j=1}^8\gamma_j=|\det A_{3,s}|^{-2}$.
Define the actual source-to-target map and its inverse by
\[
 \mathcal B_C=\mathfrak F_W(M\otimes I_8)\mathfrak F_U^{-1}
             =\mathfrak T\,\mathcal M_U,\qquad
 \mathcal B_C^{-1}=\mathfrak F_U(M^{-1}\otimes I_8)
                  \mathfrak F_W^{-1}
             =\mathcal M_U^{-1}\mathfrak T^{-1}.
 \tag{ET32}
\]
Their products are the identity on their full original domains.
The inverse metric square in coefficient coordinates is
\[
 ((M^{-1})^*M^{-1})\otimes
       (\Gamma_W^{-1}\Gamma_U),\qquad
 \Gamma_b=\diag(G_{b,r},G_{b,r}) .
 \tag{ET33}
\]
For a map from metric $I_3\otimes\Gamma_W$ to metric
$I_3\otimes\Gamma_U$, this follows by the definition of the metric
adjoint. The second factor is positive in its $\Gamma_W$ inner product;
its eigenvalues are exactly $\gamma_j^2$ by GD14.
Diagonalizing both positive metric squares proves that the complete
finite-$C$ inverse singular list is
\[
 \{s_i(C)\gamma_j:1\le i\le3,\ 1\le j\le8\}.
 \tag{ET34}
\]
Here the $s_i$ are the explicitly evaluated ET28a square roots.
Sorting these twenty-four positive numbers and multiplying the largest
$k$ gives every exact inverse exterior norm, with all original conductor
weights present, even where the order within the list changes.

Put $\widehat\gamma_{2j-1}=\widehat\gamma_{2j}=\gamma_j$ for
$1\le j\le8$. For $C\to0$, fixed nonzero $A$, fixed $B$ and the
unchanged conductor, all eight large factors eventually exceed
all sixteen small factors, because their scale ratio is $|C|^{-4}$
times a positive fixed constant. Products in ET34 therefore give
\[
 \left\|\bigwedge^k\mathcal B_C^{-1}\right\|\sim
 \begin{cases}
 d_A^k\left(\prod_{j=1}^k\gamma_j\right)|C|^{-3k},
                    &1\le k\le8,\\[3pt]
 d_A^8\,5^{-(k-8)}
 \left(\prod_{j=1}^8\gamma_j\right)
 \left(\prod_{j=1}^{k-8}\widehat\gamma_j\right)|C|^{k-32},
                    &8\le k\le24 .
 \end{cases}
 \tag{ET35}
\]
Repeated $\gamma_j$ do not affect this conclusion: ordered positive
finite lists are continuous and tied limiting products coincide.
In full degree there is the exact, not merely limiting, equality
\[
 \left\|\bigwedge^{24}\mathcal B_C^{-1}\right\|
 =\left(\frac{5|A|}{4|C|}\right)^8|\det A_{3,s}|^{-6},
 \quad
 |\det A_{3,s}|=
 \left|\frac{(-i)^v\mu_v}{v!}\right|^4
 \prod_{j=0}^3\frac{\rho_{v+j}^{(s)}}{\rho_j^{(s)}} .
 \tag{ET36}
\]
This is the actual inverse-exterior comparison with audit equation (11).
The cubic coefficient-lattice loss multiplies the retained conductor
cost; it does not set the nonzero original moment $\mu_v$ to zero.

There is a uniform finite-$C$ estimate with these exact leading scales.
Put
\[
 a_C=\frac{125|A|}{4|C|^3},\quad b_C=\frac{|C|}{5},\quad
 \delta_C=\frac{|C|^2}{25|A|^2}
       +\frac{|B|^2|C|^4}{625|A|^2}
       +\frac{16|C|^8}{390625|A|^2},\quad
 e_k=\min(k,8,24-k).
 \tag{ET37}
\]
The independent derivation ETR21--22 gives the same expression;
here is its full cross-receiver consequence.
Let $L_k$ be the product of the largest $k$ entries of the exact
baseline multiset
$\{a_C\gamma_j,b_C\gamma_j,b_C\gamma_j:1\le j\le8\}$.
Then, for every $A,C\ne0$ and every $B$,
\[
 L_k\le\left\|\bigwedge^k\mathcal B_C^{-1}\right\|
          \le L_k(1+\delta_C)^{e_k/2},
 \qquad 0\le k\le24 .
 \tag{ET38}
\]
To prove this, write $\theta=s_1/a_C$.
The largest eigenvalue of $uu^*+|y|^2\diag(1,1,0)$, where
$u=(x,z,w)^{\mathsf T}$, lies between $|w|^2$ and
$\|u\|^2+|y|^2=|w|^2(1+\delta_C)$.
Thus $1\le\theta\le\sqrt{1+\delta_C}$, and the exact finite
singular triple is
\[
 (s_1,s_2,s_3)=(a_C\theta,b_C,b_C/\theta).
 \tag{ET39}
\]
Its first entry exceeds both $a_C$ and $b_C$ by ET28a; its third
is smaller than both. Its total product is $a_Cb_C^2$.
Consequently, in each fixed $j$ block, the product of the largest
zero, one, two or three actual entries is at least the corresponding
baseline product. For one entry this uses $s_1\ge\max(a_C,b_C)$.
For two entries, $s_1b_C$ is at least $a_Cb_C$ and $b_C^2$.
For zero or three entries there is equality.
Take a subset attaining the largest baseline $k$-product, count its
entries in each block, and replace them by that many largest actual
entries in the same block. This proves the lower bound.
For the upper bound, each actual entry is its specified baseline
entry times $\theta$, $1$, or $\theta^{-1}$.
Across the eight blocks the exponent list has eight $+1$'s,
eight zeros and eight $-1$'s. A $k$-element subset has exponent
at most $\min(k,8,24-k)$, the last bound following by taking its
complement. Its baseline product is at most $L_k$.
Maximizing over all subsets proves the upper bound in ET38.
This proof holds even when baseline values change order and when
original conductor singular values coincide. At $k=24$ the two
bounds agree with the exact determinant ET36.
